\scanpage{38}
% archive sheet number: 37
\begingroup\makeatletter\tagsleft@true\makeatother
Mais si cette condition n'est pas satisfaite pour $\unsure{\deleted{\Sigma|E}}\unsure{\Pi|\pi}_1,\Pi_2$ non panachées, on reste en plan, et on regarde l'ens. à quatre éléments
\[
C=\Pi_1\amalg\Pi_2=\omega_1\amalg\omega_2
=\{\varepsilon_1^+,\varepsilon_1^-,\varepsilon_2^+,\varepsilon_2^-\}
\tag*{(17)}
\]
et la relation d'ordre dessus, où les seules relations $<$ sont données par
\[
\varepsilon_1^+<\varepsilon_2^-,\qquad\varepsilon_2^+<\varepsilon_1^-
\tag*{(22)}
\]
\diagram{\begin{tikzcd}[ampersand replacement=\&,row sep=2.3em,column sep=2.6em]
|[label=left:{\(\varepsilon_1^-\)}]| \bullet \arrow[from=2-2,to=1-1]
 \& |[label=right:{\(\varepsilon_2^-\)}]| \bullet\\
|[label=left:{\(\varepsilon_1^+\)},label=below:{\(\deleted{\varepsilon_2^+}\)}]| \bullet \arrow[ur,crossing over]
 \& |[label=right:{\(\varepsilon_2^+\)}]| \bullet
\end{tikzcd}}

1.3. Soit \deleted{$(\Pi_\alpha)_{\alpha\in A}$ une famille}
\[
(\Pi_\alpha)_{\alpha\in A},\qquad \Pi_\alpha=\{\Sigma_\alpha^\varepsilon\}_{\varepsilon\in\omega_\alpha}
\tag*{(23)}
\]
\margin{c'est un élément de $\unsure{\mathfrak{P}|\mathcal{P}}(\mathfrak{P}_2(\unsure{\deleted{X|\times}}\unsure{\mathfrak{P}|\mathcal{P}|\square})(X))$}
\unsure{une|un} \deleted{famille} \inserted{ens. (\unsure{écrit|écrite|noté} \unsure{indiciellement|individuellement})} de décompositions binaires de $X$, on suppose \deleted{qu'elles}
\[
\forall\alpha,\beta\in A,\quad\alpha\ne\beta,\qquad
\Pi_\alpha,\Pi_\beta\ \text{sont \unsure{mut|mul}\textsuperscript{t}. non panachées\unsure{.}}
\tag*{(24)}
\]
\margin{On dit, si $\operatorname{Card} A=n$, que l'on a une \uline{décomposition} \uline{$n$-aire} non panachée de $X$}

\inserted{\uline{Déf}} On dit \deleted{alors} que $(\Pi_\alpha)_{\alpha\in A}$ est non \unsure{panachée|panachées}. (Terminologie analogue pour \uline{strictement} non panachée).

Considérons
\[
C=\coprod_{\alpha\in A}\omega_\alpha
\tag*{(25)}
\]
\inserted{avec l'involution ss pt fixe $\sigma$ qui, sur chaque $\omega_\alpha$, induit une \deleted{\illegible{1}} transposition}
et sur $C$ la relation d'ordre, dans laquelle on a
\endgroup
